\documentclass[10pt,a4paper]{amsart}
\usepackage[margin=19mm]{geometry}
\usepackage{amsmath,amssymb,mathtools}
\usepackage[scr=rsfs,cal=euler]{mathalfa}
\usepackage{xfrac}
\usepackage[unicode,hidelinks,bookmarks=false]{hyperref}
\newcommand{\ie}{i.e.\ }
\newcommand{\cf}{cf.\ }
\newcommand{\resp}{respectively\ }
\newcommand{\Subsection}{\subsection}
\providecommand{\nobreakdash}{\nobreak\hbox{-}}
\setlength{\emergencystretch}{2em}
\multlinegap0pt
\usepackage{bm}
\usepackage{bbm}
\usepackage{dsfont}
\usepackage{xspace}
\usepackage{cancel}
\usepackage{mathtools}
\usepackage{amsthm}
\makeatletter
\@ifundefined{theorem}{\newtheorem{theorem}{Theorem}}{}
\@ifundefined{lemma}{\newtheorem{lemma}[theorem]{Lemma}}{}
\@ifundefined{proposition}{\newtheorem{proposition}[theorem]{Proposition}}{}
\makeatother
\makeatletter
\expandafter\ifx\csname revision\endcsname\relax
\newenvironment{revision}{\color{blue}}{\color{black}}
\fi
\makeatother
\title[Existence and Regularity of Stable Resonant Spectral Submani]{Existence and Regularity of Stable Resonant Spectral Submanifolds and Linearization Maps}
\author{Florian Kogelbauer, Rafael de la Llave}
\date{Source: arXiv:2609.18236v1 (CC BY 4.0)}
\begin{document}
\maketitle

\noindent\emph{Source and licence.} Existence and Regularity of Stable Resonant Spectral Submanifolds and Linearization Maps. Version arXiv:2609.18236v1 (CC BY 4.0), distributed under the Creative Commons Attribution 4.0 International licence, as recorded in the arXiv licence metadata for this exact version and shown on the abstract page https://arxiv.org/abs/2609.18236v1. Full rights notice: licenses/arxiv-2609.18236-CC-BY-4.0.txt.\par\smallskip

\noindent\textit{Editorial scope.} Counted unit: the Proposition whose source label is \texttt{propfixedpoint} in the article (source0.tex lines 1417--1423, source label \texttt{propfixedpoint}), delivered together with the article's own complete original proof of it (source0.tex lines 1425--1584, one \texttt{proof} environment of 160 lines). No mathematics is written, repaired or re-derived: statement and proof are the article's own text line for line. Required context retained uncounted: mainthm (lines 88--112); asslambda (lines 95--97); lemma\_finite\_resonances (lines 229--231); Lemma\_SM (lines 1119--1150); coholdelta (lines 1273--1275); N\_f\_expand (lines 1281--1283); lemma\_affine\_tail (lines 1343--1373); N\_delta (lines 1353--1359); bound\_difference\_N\_delta (lines 1367--1372). Every definition, display and label of those lines is retained unchanged. Omitted from this package and not redistributed: 1--87, 98--228, 232--1118, 1151--1272, 1276--1280, 1284--1342, 1360--1366, 1373--1416, 1424--1424, 1585--2207. Nothing inside a retained excerpt is omitted or altered: every retained body is delivered exactly as the article's lines are written, so no omission receipt of any kind accompanies this package. Citations: the retained excerpts carry no citation command at all, so no bibliography entry of the article is transcribed and no reference is redistributed. Reference adaptation: none was needed and none was made; every reference inside the delivered unit resolves inside it, so no unresolved reference or citation is produced. Printed slips kept literal: no printed word, symbol, hypothesis or number of the counted statement or of its proof was altered. Layout and class: the article's own class is \texttt{amsart} with packages including comment, xcolor, amssymb, graphicx; the wrapper is the lane's pinned \texttt{amsart} editorial wrapper at 10pt on A4, which restates the theorem-like environments the retained text needs and adds the article's own macro definitions that the retained text needs (none), with extra wrapper packages (bm, bbm, dsfont, xspace, cancel, mathtools). The delivered numbering is the unit's own. Assets: no retained span contains a figure environment, a graphics-inclusion command, a PDF-page inclusion, a table or a TikZ picture; no image file is copied into this package, the package declares no asset dependency and no image file is redistributed with it. Licence: version arXiv:2609.18236v1 of this article is distributed under the Creative Commons Attribution 4.0 International licence, as recorded in the arXiv licence metadata for this exact version and shown on the abstract page https://arxiv.org/abs/2609.18236v1. A full rights notice is delivered at licenses/arxiv-2609.18236-CC-BY-4.0.txt. Characters: every retained excerpt is pure ASCII, so the delivered engine, XeLaTeX with its default Latin Modern text font, reports no missing character.
\begin{theorem}\label{mainthm}
Let $f:\mathbb{C}^s \to\mathbb{C}^s$ be an analytic map with $f(0)=0$, understood, when $f$ represents a real system, in complex eigencoordinates compatible with the underlying real structure. All differentiability assertions below are in the real sense. Let
\begin{equation}
A=Df(0),
\end{equation}
be semisimple and invertible.
Let $X_1$ be an $A$-invariant, $d$-dimensional subspace spanned by eigenvectors of $A$ associated to eigenvalues $\Lambda = \{\lambda_1,...,\lambda_d\}$ such that
\begin{equation}\label{asslambda}
|\lambda_j|< 1,\quad 1\leq j\leq d.
\end{equation}
Assume first that $I_{\rm res}(\Lambda)\neq\varnothing$. For each resonant multi-index $n\in I_{\rm res}(\Lambda)$, define
\begin{equation}\label{def_r_star}
    r(n):=
\min\left\{
|n|-1,\,
\min_{j\in\operatorname{supp}(n)} n_j
\right\} \geq 1,\quad r_*:=
\min_{n\in I_{\rm res}(\Lambda)} r(n).
\end{equation}
Then, on a sufficiently small neighborhood of the origin in $X_1$, there exists a parametrization $h$ of class $C^{r_*,1-\varepsilon}$ for every $\varepsilon\in(0,1)$ satisfying
\begin{equation}
        h(0)=0,\qquad Dh(0)=\iota_{X_1},\qquad f\circ h=h\circ A|_{X_1},
\end{equation}
where $\iota_{X_1}$ denotes the canonical inclusion embedding of $X_1$ into the ambient space. After restricting the domain, $h$ is an embedding and its image is an $f$-invariant manifold tangent to $X_1$. The parametrization admits logarithmic-polynomial approximations to arbitrary finite order. The stated regularity is a uniform lower bound on the regularity of this invariance-linearization parametrization. If $I_{\mathrm{res}}(\Lambda)=\varnothing$, the parametrization can be chosen analytic.
\end{theorem}

\begin{lemma}\label{lemma_finite_resonances}
Under \eqref{asslambda} and the invertibility of $A$, the set $I_{\rm res}(\Lambda)$ is finite.
\end{lemma}

\begin{proposition}\label{Lemma_SM}
Let $A\in\mathbb C^{s\times s}$ be semisimple and invertible and let $\lambda_1,\ldots,\lambda_s$ denote the eigenvalues of $A$, counted with
multiplicity. Let
$X_1\subset\mathbb C^s$ be a $d$-dimensional $A$-invariant subspace associated with
eigenvalues $\lambda_1,\ldots,\lambda_d$, and assume
\begin{equation}
        |\lambda_i|<1,
        \qquad
        1\leq i\leq d.
\end{equation}
Suppose $N$ is chosen so large that
\begin{equation}\label{nonres_tail}
        \lambda^n\ne \lambda_j,
        \qquad
        |n|\geq N,
        \qquad
        1\leq j\leq s.
\end{equation}
Then, for $\beta>0$ sufficiently small, the operator
\begin{equation}
        \mathcal S_A:\Gamma_N^d\to\Gamma_N^d,
        \qquad
        \mathcal S_A H(x)=AH(x)-H(Ax),
\end{equation}
is bounded and invertible. Moreover, there exists $q_*\in(0,1)$ such that
\begin{equation}\label{SA_inverse_bound}
        \|\mathcal S_A^{-1}\|_{\rm op}
        \le
        \frac{1}{1-q_*}
        \frac{1}{\min_{1\leq j\leq s}|\lambda_j|}.
\end{equation}
\end{proposition}

\begin{equation}\label{coholdelta}
Ah^\delta(x)+N_f(h^\delta(x))=h^\delta(Ax).
\end{equation}

\begin{equation}\label{N_f_expand}
N_f^\delta(x)=\sum_{|n|=2}^\infty \delta^{|n|-2} f_nx^n,
\end{equation}

\begin{lemma}[Affine tail composition]\label{lemma_affine_tail}
Let $h^{<N}$ be a fixed logarithmic polynomial with
\begin{equation}
        h^{<N}(x)=Lx+O(|x|^2),
\end{equation}
and let $H\in\Gamma_N^d$. Let $N_f^\delta$ be analytic near the origin and satisfy
\begin{equation}
        N_f^\delta(z)=O(|z|^2).
\end{equation}
Then
\begin{equation}\label{N_delta}
        \mathcal N_\delta(H)
        =
        -\delta\left[
        N_f^\delta(h^{<N}+H)-N_f^\delta(h^{<N})
        \right]
\end{equation}
belongs to $\Gamma_N^d$. Moreover, for every $R>0$ sufficiently small, there exists
$C_R>0$ such that, for all $H_1,H_2\in\Gamma_N^d$ with
\begin{equation}
        \|H_i\|_{\Gamma_N^d,\beta,\sigma}\leq R,
        \qquad i=1,2,
\end{equation}
one has
\begin{equation}\label{bound_difference_N_delta}
        \|\mathcal N_\delta(H_1)-\mathcal N_\delta(H_2)\|_{\Gamma_N^d,\beta,\sigma}
        \le
        \delta C_R
        \|H_1-H_2\|_{\Gamma_N^d,\beta,\sigma}.
\end{equation}
\end{lemma}

\begin{proposition}\label{propfixedpoint}
Let the assumptions of Theorem \ref{mainthm} and Proposition \ref{Lemma_SM} with
$\beta$ sufficiently small be met. Assume further that the invariance equation has been
scaled according to \eqref{coholdelta} for $\delta>0$ sufficiently small. Then, for any
approximate solution $h^{<N}$, with $N\in\mathbb N$ chosen larger than the order of every resonant index, which is possible by Lemma~\ref{lemma_finite_resonances}) and sufficiently large for Proposition~\ref{Lemma_SM}, the functional $\mathcal T_\delta(\cdot,h^{<N})$ is a contraction from the
closed unit ball $B^1(\Gamma_N^d)$ into itself and thus has a unique fixed point.
\end{proposition}

\begin{proof}
To avoid cluttering the notation, we write
\begin{equation}
        \|\cdot\|=\|\cdot\|_{\Gamma_N^d,\beta,\sigma}.
\end{equation}
By Lemma \ref{lemma_affine_tail}, $\mathcal N_\delta(H)\in\Gamma_N^d$ for
$H\in\Gamma_N^d$, and $\tau_\delta\in\Gamma_N^d$ because $h^{<N}$ solves the scaled
invariance equation up to order $N-1$. Thus $\mathcal T_\delta$ is well-defined as a map
from $\Gamma_N^d$ to itself.

We first prove that $\mathcal T_\delta$ is a contraction on $B^1(\Gamma_N^d)$. Let
$H,\widetilde H\in B^1(\Gamma_N^d)$. By Lemma \ref{lemma_affine_tail}, applied with
$R=1$, there exists a constant $C_1>0$ such that
\begin{equation}\label{contraction_revised}
\begin{split}
        \|\mathcal T_\delta(H,h^{<N})-\mathcal T_\delta(\widetilde H,h^{<N})\|
        &=
        \left\|
        \mathcal S_A^{-1}
        \left(
        \mathcal N_\delta(H)-\mathcal N_\delta(\widetilde H)
        \right)
        \right\| \\
        &\le
        \|\mathcal S_A^{-1}\|_{\rm op}
        \|\mathcal N_\delta(H)-\mathcal N_\delta(\widetilde H)\| \\
        &\le
        \delta
        \|\mathcal S_A^{-1}\|_{\rm op}
        C_1
        \|H-\widetilde H\|.
\end{split}
\end{equation}
Choosing $\delta>0$ sufficiently small so that
\begin{equation}
        \delta
        \|\mathcal S_A^{-1}\|_{\rm op}
        C_1
        <1,
\end{equation}
we obtain that $\mathcal T_\delta(\cdot,h^{<N})$ is a contraction on
$B^1(\Gamma_N^d)$.\\
It remains to show that $\mathcal T_\delta$ maps the closed unit ball into itself. For
$H\in B^1(\Gamma_N^d)$, we have
\begin{equation}\label{ball_revised}
        \mathcal T_\delta(H,h^{<N})
        =
        \mathcal S_A^{-1}\tau_\delta
        +
        \mathcal S_A^{-1}\mathcal N_\delta(H).
\end{equation}
The approximate solution $h^{<N}$ is constructed for the scaled
invariance equation through order $N-1$. Therefore the residual $\tau_\delta$ starts at order $N$, and hence
\begin{equation}
    \tau_\delta\in\Gamma_N^d.
\end{equation}
Moreover, by \eqref{N_f_expand}, the homogeneous term of order
$k$ in the scaled nonlinear part $\delta N_f^\delta$ carries the
factor $\delta^{k-1}$. For completeness, we make the dependence on the scaling parameter $\delta$ explicit. Writing again
\begin{equation}
N_f^\delta(z)
=
\sum_{k\geq 2} B_k^\delta(z,\ldots,z),
\qquad
B_k^\delta=\delta^{k-2}B_k,
\end{equation}
we have, for $0<\delta\leq 1$, the uniform bounds
\begin{equation}
\|B_k^\delta\|\leq \|B_k\|.
\end{equation}
Since $h^{<N}$ is a fixed logarithmic polynomial, multiplication by
$h^{<N}$ defines a bounded operator on $\Gamma_N^d$. Hence, after
restricting the domain if necessary, analyticity of $N_f$ implies that,
for $\|H_i\|_{\Gamma_N^d}\leq R$,
\begin{equation}
\bigl\|
N_f^\delta(h^{<N}+H_1)
-
N_f^\delta(h^{<N}+H_2)
\bigr\|_{\Gamma_N^d}
\leq
C_R\|H_1-H_2\|_{\Gamma_N^d},
\end{equation}
where $C_R$ can be chosen independently of sufficiently small $\delta$.
Multiplication by the prefactor $\delta$ in \eqref{N_delta} therefore gives
\eqref{bound_difference_N_delta}. Moreover, the order-by-order construction of $h^{<N}$ shows inductively that its homogeneous coefficient of order $k$ is of
the form
\begin{equation}
h_k^\delta=\delta^{k-1}\widehat h_k,
\end{equation}
with $\widehat h_k$ independent of $\delta$. Indeed, an order-$k$
contribution arising from an $m$-linear term of $N_f$, evaluated on
homogeneous terms of orders $j_1,\ldots,j_m$ with
$j_1+\cdots+j_m=k$, carries the factor
\begin{equation}
\delta^{m-1}\prod_{i=1}^m\delta^{j_i-1}
=
\delta^{k-1},
\end{equation}
and inversion of the corresponding cohomology operator does not alter
this factor. Since all terms of orders smaller than $N$ cancel by
construction, the residual $\tau_\delta$ therefore starts at order $N$,
and its coefficients of order $k\geq N$ carry a factor $\delta^{k-1}$. The same analytic majorant estimate as above then yields, uniformly for sufficiently small $\delta$,
\begin{equation}\label{bound_residual_revised}
\|S_A^{-1}\tau_\delta\|_{\Gamma_N^d}
\leq
C_\tau\delta^{N-1},
\end{equation}
for some constant $C_\tau>0$. For the nonlinear tail, Lemma \ref{lemma_affine_tail} gives, again with $R=1$,
\begin{equation}
        \|\mathcal N_\delta(H)-\mathcal N_\delta(0)\|
        \le
        \delta C_1\|H\|.
\end{equation}
Since $\mathcal N_\delta(0)=0$, this yields
\begin{equation}\label{tail_bound_revised}
        \|\mathcal S_A^{-1}\mathcal N_\delta(H)\|
        \le
        \delta
        \|\mathcal S_A^{-1}\|_{\rm op}
        C_1
        \|H\|.
\end{equation}
Combining \eqref{ball_revised}, \eqref{bound_residual_revised}, and
\eqref{tail_bound_revised}, we obtain
\begin{equation}
    \|T_\delta(H,h^{<N})\|
    \leq
    C_\tau\delta^{N-1}
    +
    \delta\|S_A^{-1}\|_{\rm op}C_1\|H\|.
\end{equation}
For $H\in B^1(\Gamma_N^d)$, this gives
\begin{equation}
        \|\mathcal T_\delta(H,h^{<N})\|
        \le
        C_{\tau}\delta^{N-1}
        +
        \delta
        \|\mathcal S_A^{-1}\|_{\rm op}
        C_1 .
\end{equation}
By decreasing $\delta>0$ further, if necessary, we can ensure that
\begin{equation}
        C_{\tau}\delta^{N-1}
        +
        \delta
        \|\mathcal S_A^{-1}\|_{\rm op}
        C_1
        \leq 1.
\end{equation}
Thus $\mathcal T_\delta(\cdot,h^{<N})$ maps the closed unit ball
$B^1(\Gamma_N^d)$ into itself.\\
The Banach fixed point theorem now gives a unique fixed point
$H\in B^1(\Gamma_N^d)$. Consequently,
\begin{equation}
        h=h^{<N}+H
\end{equation}
solves the scaled invariance equation. This proves the claim.
\end{proof}

\end{document}
